\section{La filosofía de la simbiosis entre el ser humano y la máquina}

\subsection{La historia de la evolución de las herramientas es la historia de la humanidad}

Yang Shuo (杨硕) planteó un punto de vista filosófico profundo: \textbf{el desarrollo de la civilización humana es, en esencia, la historia de la evolución de las herramientas (las máquinas)}.

Describió una cadena evolutiva que va de las herramientas primitivas a los robots modernos:
\begin{itemize}
    \item Lascas de piedra, palos $\to$ el hacha (atar el palo y la piedra juntos)
    \item Ruedas de madera, aperos de hierro $\to$ casas, relojes
    \item Máquinas de escribir, impresoras $\to$ automóviles, trenes, aviones
    \item Refrigeradores, lavadoras $\to$ computadoras, teléfonos inteligentes $\to$ robots
\end{itemize}

«¿Cómo se dice que la historia se desarrolla cada vez mejor? En realidad no es que el ser humano mismo experimente algún tipo de evolución, sino que las máquinas que el ser humano fabrica son cada vez mejores.» Por lo tanto, la mejora en la capacidad y la inteligencia de los robots representa precisamente la mejora en la capacidad y la inteligencia de la especie humana.

\begin{importantbox}{La civilización de las máquinas es la civilización humana}
La tesis filosófica central de Yang Shuo: la civilización humana ha estado ligada a la civilización de las máquinas desde el principio; ambas son una sola cosa. Antes de que existieran las herramientas, el ser humano no se distinguía de los monos; fue la aparición de las herramientas lo que separó al ser humano de los animales. Por lo tanto, desarrollar mejores robots no es una amenaza para la humanidad, sino una manifestación necesaria del progreso de la civilización humana. «Lo que existe no es la historia de la humanidad, sino la historia del desarrollo de las máquinas.»
\end{importantbox}

\subsection{La desencarnación del ser humano y la antropomorfización de la máquina}

Yang Shuo planteó una tendencia que invita a la reflexión: el ser humano se está «desencarnando» de manera continua, mientras que la máquina se está «antropomorfizando» de manera continua, y al final ambas convergerán en algún punto.

Puso un ejemplo muy vívido: si te hospedas en un hotel y encuentras un cabello en la tina, sientes un rechazo instintivo (es la repulsión instintiva del ser humano hacia los desechos corporales de otras personas); pero si lo que encuentras es un trozo de cable, no lo sientes sucio. Si el huésped anterior hubiera sido un robot en lugar de un ser humano, incluso sentirías que el lugar está limpio.

Esto muestra que en lo más profundo de la naturaleza humana hay una búsqueda por «eliminar» el propio cuerpo y por lograr una mejor conexión con la máquina. Yang Shuo especuló que, dentro de algunos miles de años, la mayor parte de los «seres humanos» que caminen sobre la Tierra podrían ser mecanizados o semimecanizados: la conciencia se habrá trasladado a un cuerpo mecánico, pero la civilización humana seguirá existiendo.

Resumió el futuro de los robots con tres palabras clave: \textbf{simbiosis} (el ser humano y la máquina coexisten y dependen uno del otro en distintos niveles) y \textbf{muerte} (si la vida digital necesita que se le otorgue la muerte para mantener su capacidad de creación). Al preguntársele «¿el ser humano se vuelve cada vez más desencarnado y la máquina cada vez más parecida al ser humano?», Yang Shuo afirmó «sí», porque estas dos cosas «en esencia son una sola cosa».

\subsection{¿Necesita morir la máquina?}

Esta es la discusión de mayor profundidad filosófica de toda la entrevista. Yang Shuo planteó la segunda de las tres palabras clave sobre el futuro de los robots: \textbf{la muerte}.

\begin{warningbox}{¿Necesita otorgársele la muerte a la vida digital?}
Yang Shuo considera que la futura sociedad dominada por máquinas quizá necesite discutir una pregunta central: ¿se debe dejar morir a la vida digital? La conciencia de los robots puede copiarse sin cesar y migrar entre distintos servidores; en esencia, es inmortal. Pero la sociedad de las máquinas podría descubrir que, para alcanzar la creatividad al nivel del ser humano, \textbf{es necesario que exista de manera constante la amenaza de la muerte}. Una vida sin límite de tiempo siempre puede «dejarlo para mañana» y carece de un impulso creativo urgente. Esto se hace eco de la célebre frase de Steve Jobs: «La muerte es el mejor invento de la vida.» Este problema no necesita resolverse en la actualidad, pero a largo plazo podría ser uno de los mayores problemas sociales que enfrente la civilización de las máquinas.
\end{warningbox}

Yang Shuo predice que la sociedad de las máquinas quizá necesite fijar un límite de tiempo de 50 u 80 años para cada vida mecánica, al término del cual sea obligatorio borrar la conciencia. Pero esto necesariamente generará una enorme controversia: un robot inmortal no aceptará voluntariamente su propia muerte, y cómo llevar a cabo esta «muerte digital» será un problema social de gran magnitud.

\subsection{Simbiosis: la respuesta última a la relación entre el ser humano y la máquina}

Yang Shuo resumió con la palabra \textbf{simbiosis} su juicio último sobre la relación entre el ser humano y la máquina. Esto abarca varios niveles:

\begin{enumerate}
    \item \textbf{Ya existe simbiosis en el presente}: la ropa, la comida, la vivienda y el transporte del ser humano son provistos por máquinas; sin refrigerador, televisor a color ni lavadora, la persona moderna no puede vivir
    \item \textbf{La simbiosis es bidireccional}: así como el ser humano no aceptaría destruir todos los aparatos eléctricos, en el futuro los robots conscientes tampoco aceptarán eliminar a todos los seres humanos
    \item \textbf{La máquina se irá acercando al modo de vida del ser humano}: incluido, posiblemente, tener que enfrentar la muerte
    \item \textbf{El ser humano se irá acercando a la forma de la máquina}: la desencarnación, la mecanización, la semimecanización
\end{enumerate}

\subsection{La conexión emocional con los robots}

Un detalle conmovedor de la entrevista es cuando Yang Shuo mencionó cómo se sintió al dejar el Tesla Optimus: «Estuve muy mal muchos días, sentía que había convivido con este robot día y noche, y ahora tenía que dejarlo.» El presentador también compartió una experiencia similar: al hacer experimentos, desarrolló una conexión emocional con el robot, «aunque fuera muy torpe y no entendiera lo que le decía».

Yang Shuo considera que este tipo de conexión emocional entre el ser humano y el robot ocurrirá de manera generalizada en el futuro, tal como la despedida es inevitable entre los seres humanos.

\subsection{Resumen de la sección}

La reflexión filosófica de Yang Shuo sobre la relación entre el ser humano y la máquina trasciende el plano técnico y toca las cuestiones centrales del existencialismo. La historia de la evolución de las herramientas como historia de la humanidad, la tendencia a la desencarnación del ser humano y la antropomorfización de la máquina, la pregunta de si la vida digital necesita morir, y la «simbiosis» como respuesta última a la relación entre el ser humano y la máquina: estas reflexiones ofrecen a la industria de la robótica una perspectiva de profundidad que va mucho más allá de la dimensión técnica.
